\section{Note on sources and delivery}
This finished homily addresses an adult parish assembly for All Saints, 1 November 2026, under the U.S. \work{Roman Missal}, Third Edition. Its spoken body contains \textbf{1,174 words}, counted as whitespace-separated words, excluding the title and this terminal apparatus. At an estimated unhurried 105--115 words per minute, delivery would take approximately \textbf{10.2--11.2 minutes}, before additional pauses. The author reviewed the prose silently for sense, oral sequence and manageable sentences. No audible rehearsal or timed human delivery is claimed.

The argument derives from the reviewed expansive study's ``Purified desire and the sight of God'' and the concise study's comparisons of purification, vision and mercy. It follows Augustine's explanation of 1 John 3:1--3: real present sonship, future vision and willing purification through God's help. The decisive detail is that John attributes purification to the person who has already received adoption. This changes holiness from a distant achievement contemplated in the saints into present cooperation with the grace that brought them to God. The application of giving needed help despite resentment derives from Augustine's explanation of Matthew 5:7. The speech's connections among Revelation, the psalm, First John, the Beatitudes and the Mass's prayers are the homily's own theological synthesis over the reviewed evidence. No Father is attributed an interpretation of the complete modern Mass.

The source and rights qualifications of the canonical study remain applicable. Scriptural phrases follow the public-domain Douay--Rheims/Challoner study witness, rather than reproduce the approved U.S. Lectionary. No complete prayer body or substitute translation is supplied. The prayer descriptions retain the reviewed limits of the qualified secondary Missal witness; they do not claim a new authentication of every word of the current altar book. Historical NPNF and Marshall translations remain outside the project license.

\section{References}
\begin{itemize}[leftmargin=1.3em,itemsep=.4em]
\item \textbf{Scripture.} \work{Douay--Rheims Bible}, Challoner text, Project Gutenberg eBook 1581, retained verse-text edition: Revelation 7:2--4,9--14; Psalm 23:1--6 (Hebrew Psalm 24); 1 John 3:1--3; Matthew 5:1--12a. The U.S. appointments and partial-verse boundaries are identified in the canonical study; the speech supplies commentary rather than full appointed readings.
\item \textbf{Augustine.} \work{Tractates on the First Epistle of John}, IV.4--7, especially IV.6 on enlarged desire and the vessel, and IV.7 on God's help and willing purification, trans. H. Browne, NPNF, first series, VII. \work{Our Lord's Sermon on the Mount}, I.2.7--8 and I.3.10, on mercy, purity and helping the weaker while needing help, trans. William Findlay, NPNF, first series, VI. Retained historical CCEL English deliveries; no new original-language collation.
\item \textbf{Bede.} \work{Explanation of the Apocalypse}, trans. Edward Marshall (Oxford and London: James Parker, 1878), I, complete comment on 7:14, p. 53: the Lamb's blood cleanses the whole Church, beyond martyrs alone. Checked retained transcription; no new original-Latin collation.
\item \textbf{Liturgical prayers.} U.S. \work{Roman Missal}, Third Edition (2011), All Saints: Collect and Prayer after Communion. Qualified secondary Latin witness: 2002 reproduction, PDF pp. 512--514, read with the official 2008 \work{Variationes}; full evidence, missing-exemplar qualifications and exact locators remain with the canonical study and its same-edition shared Missal owner. Prayer subjects are explained without reconstructing protected wording.
\end{itemize}
